\documentclass[12pt]{article}
\usepackage[utf8]{inputenc}
\usepackage[T1]{fontenc}
\usepackage{hyperref}
\title{Yaleba Language Resources for K-12 Curriculum Development}
\author{Compiled by Claude Sonnet 5.5 (claude-sonnet-5-5) for Andr\'as Kornai}
\date{}
\begin{document}
\maketitle

\section{Introduction}
Yaleba (ISO 639-3: ylb) is an Austronesian language of Milne Bay Province, Papua New Guinea. Joshua Project gives the alternate name Aisana and the dialects Gamadoudou, Sileba and Soma'a, and lists about 2,400 people. Wikipedia's Wagawaga article gives 1,900 native speakers, citing the 2000 census, but that figure covers the broader Wagawaga entry. The two figures are inconsistent. The Yaleba live on a peninsula south of Alotau, and most are Christian. The search found very little material specific to this language.

\section{Primary Educational Resources}
No primers, textbooks or school materials in Yaleba turned up in the search. The reference description says the language is not known to be taught in schools.

\section{Linguistic and Reference Materials}
No dictionary, grammar or orthography for Yaleba turned up in the search. The only reference-type pages found are the code disambiguation page and the related Wagawaga article, which do not describe Yaleba in detail:
\begin{itemize}
\item Wikipedia, ``YLB'' (code disambiguation): \url{https://en.wikipedia.org/wiki/YLB}
\item Wikipedia, ``Wagawaga language (New Guinea)'': \url{https://en.wikipedia.org/wiki/Wagawaga_language_(New_Guinea)}
\end{itemize}

\section{Religious and Community Resources}
Joshua Project lists the scripture status as translation started, with no Audio Bible or Jesus Film available. It says God's word still needs to be translated into Yaleba.
\begin{itemize}
\item Joshua Project, Yaleba in Papua New Guinea: \url{https://joshuaproject.net/people_groups/20548}
\item Joshua Project, language page for ylb: \url{https://joshuaproject.net/languages/ylb}
\end{itemize}

\section{Digital and Audio Resources}
No audio, video, corpus or software resources for Yaleba turned up in the search.

\section{Curriculum Development Context}
Very little is documented about Yaleba. Any curriculum work would need new documentation of the language first. That would cover orthography, a word list and recorded texts, and it should be done with the community. The resources above are demographic and religious summaries, not teaching materials. The Wagawaga relationship and the speaker numbers should be checked before any planning relies on them.

\end{document}